% =============================================================================
% CHAPITRE 1 — CONTEXTE ET OBJECTIFS DE L'ALTERNANCE
% Première partie : contexte général et présentation du sujet
% =============================================================================

\chapter{Contexte et objectifs}
\label{chap:contexte_objectifs}

Ce premier chapitre présente le cadre général dans lequel s'est déroulée
l'alternance. Il introduit tout d'abord l'organisme d'accueil, son
positionnement dans le domaine de la géomatique ainsi que l'environnement
professionnel et technique de la mission. Il présente ensuite le produit Scan,
la place historique de FME dans son fonctionnement et les raisons ayant conduit
à engager une migration progressive de certains traitements vers une solution
Python autonome.

L'objectif est de situer la mission dans son contexte avant de détailler, dans
les chapitres suivants, les choix techniques effectués, les développements
réalisés ainsi que la démarche de validation mise en œuvre.

% =============================================================================
% 1.1 PRÉSENTATION DU CONTEXTE GÉNÉRAL
% =============================================================================

\section{Présentation du contexte général}
\label{sec:contexte_general}

% -----------------------------------------------------------------------------
% 1.1.1 Présentation de l'organisme d'accueil
% -----------------------------------------------------------------------------

\subsection{Présentation de l'organisme d'accueil}
\label{subsec:organisme_accueil}

Isogeo est une entreprise française spécialisée dans la gestion, la
documentation et la valorisation des données géographiques. Elle développe une
solution de géocatalogage destinée aux organisations qui produisent,
administrent ou utilisent un volume important de données géographiques.

Son activité répond à une problématique fréquemment rencontrée au sein des
systèmes d'information géographique : les données sont souvent réparties entre
plusieurs serveurs, bases de données, répertoires de fichiers et applications
métiers. Cette dispersion peut rendre difficile l'identification des ressources
disponibles, la compréhension de leur contenu, l'évaluation de leur qualité et
leur réutilisation par les différents utilisateurs.

La solution proposée par Isogeo vise ainsi à automatiser la constitution et la
mise à jour d'un catalogue de données géographiques. Elle permet aux
organisations de disposer d'un inventaire centralisé de leur patrimoine de
données, d'enrichir les fiches de métadonnées associées, puis de faciliter leur
recherche, leur consultation et leur diffusion.
\footnote{%
    \href{https://www.isogeo.com/}
    {Site officiel d'Isogeo, présentation de la solution de catalogage,
    consulté en juillet 2026.}%
}

Ce positionnement place l'entreprise à l'intersection de plusieurs domaines :

\begin{itemize}
    \item la géomatique et les systèmes d'information géographique ;
    \item la gouvernance et la qualité des données ;
    \item le développement de logiciels spécialisés ;
    \item la gestion et la publication de métadonnées ;
    \item l'intégration de solutions au sein des environnements informatiques
          des clients.
\end{itemize}

Les utilisateurs de la solution peuvent appartenir aussi bien au secteur public
qu'au secteur privé. Il peut notamment s'agir de collectivités territoriales,
d'opérateurs de réseaux, d'établissements publics, de bureaux d'études ou de
grandes organisations disposant d'un patrimoine important de données
géographiques.

\medskip

\noindent
\textbf{Principales activités}

\medskip

Les activités d'Isogeo ne se limitent pas à la mise à disposition d'une
interface de saisie de métadonnées. Elles couvrent plus largement l'ensemble du
cycle de gestion d'un catalogue de données géographiques :

\begin{itemize}
    \item l'inventaire automatisé des données présentes dans le système
          d'information d'une organisation ;
    \item la création et la mise à jour de fiches de métadonnées ;
    \item la centralisation et l'organisation des catalogues ;
    \item la recherche et la consultation des données disponibles ;
    \item la diffusion des catalogues au sein d'une organisation ou vers le
          public ;
    \item l'intégration des métadonnées dans les outils SIG et les applications
          métiers ;
    \item l'accompagnement des clients dans la structuration et la gouvernance
          de leur patrimoine de données.
\end{itemize}

Cette approche permet de relier les données effectivement présentes dans les
infrastructures techniques aux informations documentaires nécessaires à leur
compréhension et à leur exploitation.

\medskip

\noindent
\textbf{Produits et services proposés}

\medskip

L'écosystème Isogeo est constitué de plusieurs composants complémentaires.
La plateforme centrale permet d'administrer les catalogues, de consulter les
fiches de métadonnées et de gérer leur enrichissement. Le produit Scan assure
quant à lui l'inventaire des sources de données et l'extraction automatisée
d'informations techniques.

D'autres composants facilitent ensuite la consultation et la diffusion des
catalogues. Il s'agit notamment d'OpenCatalog, des plugins destinés aux logiciels
QGIS et ArcGIS Pro, ainsi que de différents modules et connecteurs permettant
d'intégrer les métadonnées Isogeo dans des applications tierces.
\footnote{%
    \href{https://help.isogeo.com/}
    {Base de connaissance officielle des produits Isogeo, consultée en
    juillet 2026.}%
}

L'entreprise propose également des prestations d'accompagnement, de
configuration, de formation et d'assistance. Ces services permettent d'adapter
le déploiement de la solution aux infrastructures, aux pratiques et aux enjeux
propres à chaque organisme.

% -----------------------------------------------------------------------------
% 1.1.2 Environnement de travail
% -----------------------------------------------------------------------------

\subsection{Environnement de travail}
\label{subsec:environnement_travail}

L'alternance s'est déroulée au sein de l'équipe Produit d'Isogeo. Cette équipe
intervient dans la conception, le développement, la maintenance et
l'amélioration continue des différents composants de la solution.

Mon travail a été réalisé sous la responsabilité de
\ReportCompanyTutor, maître d'apprentissage, et en collaboration avec les
membres de l'équipe impliqués dans le développement et la maintenance du
produit Scan.

L'intégration au sein de l'équipe Produit a permis de confronter les
développements réalisés aux contraintes d'un logiciel utilisé dans des
environnements clients réels. La mission ne consistait donc pas uniquement à
produire des scripts fonctionnels. Les solutions proposées devaient également
être robustes, maintenables, testables et compatibles avec l'architecture
existante.

\medskip

\noindent
\textbf{Organisation du travail}

\medskip

Le travail était organisé de manière progressive. Les traitements à étudier
étaient découpés en tâches distinctes en fonction des sources de données, des
fonctionnalités attendues et de leur niveau de priorité.

Pour chaque traitement, la démarche générale suivie était la suivante :

\begin{enumerate}
    \item analyser le fonctionnement du traitement FME existant ;
    \item identifier ses paramètres d'entrée et son format de sortie ;
    \item repérer les règles métier et les cas particuliers ;
    \item rechercher les bibliothèques ou fonctions Python adaptées ;
    \item développer une première implémentation ;
    \item tester le traitement sur plusieurs jeux de données ;
    \item comparer les résultats avec ceux produits par FME ;
    \item corriger les écarts et documenter la solution.
\end{enumerate}

Des échanges réguliers avec l'équipe permettaient de présenter l'avancement,
de valider les orientations techniques et d'identifier les éventuelles
incompatibilités avec le fonctionnement général de Scan.

Le suivi des développements reposait également sur un système de tickets. Chaque
ticket décrivait une fonctionnalité, une anomalie ou une amélioration attendue.
Cette organisation facilitait la priorisation des travaux et permettait de
conserver une trace des décisions prises.

\medskip

\noindent
\textbf{Outils utilisés}

\medskip

La réalisation de la mission a nécessité l'utilisation de plusieurs catégories
d'outils :

\begin{itemize}
    \item des environnements de développement Python, principalement pour
          l'écriture, le débogage et l'exécution des scripts ;
    \item Git pour la gestion des versions du code et le suivi des
          modifications ;
    \item FME Form pour consulter et exécuter les workbenches servant de
          référence ;
    \item GDAL et OGR pour la lecture et le traitement de nombreux formats de
          données géographiques ;
    \item ArcGIS Pro et ArcPy pour certains traitements propres à
          l'environnement Esri ;
    \item des systèmes de gestion de bases de données tels que PostgreSQL avec
          PostGIS, Oracle et Microsoft SQL Server ;
    \item Docker et des machines virtuelles pour reproduire certains
          environnements de test ;
    \item PyInstaller pour préparer la génération d'un exécutable Windows
          autonome.
\end{itemize}

L'utilisation conjointe de ces outils était nécessaire pour reproduire la
diversité des environnements pouvant être rencontrés lors de l'installation de
Scan chez un client.

\medskip

\noindent
\textbf{Collaboration et contrôle de la qualité}

\medskip

La collaboration au sein de l'équipe reposait sur des échanges techniques, des
revues de code et des phases de test. Les modifications étaient enregistrées
dans le gestionnaire de versions afin de pouvoir suivre leur évolution, les
comparer et les intégrer progressivement à la branche principale du projet.

Une attention particulière était portée à la lisibilité du code, à la gestion
des erreurs et à la stabilité des interfaces entre les différents composants.
Les scripts développés devaient en effet s'intégrer dans une chaîne de
traitement existante et respecter des formats d'entrée et de sortie déjà
utilisés par d'autres parties de Scan.

La documentation constituait également un élément important du travail. Elle
permettait d'expliquer le comportement des traitements, les dépendances
nécessaires, les paramètres disponibles et les limites identifiées.

% -----------------------------------------------------------------------------
% 1.1.3 Approche RSE de l'organisme d'accueil
% -----------------------------------------------------------------------------

\subsection{Approche RSE de l'organisme d'accueil}
\label{subsec:approche_rse}

La responsabilité sociétale des entreprises, généralement désignée par le sigle
RSE, correspond à la prise en compte des enjeux environnementaux, sociaux et
économiques dans les activités et les décisions d'une organisation.

La présente analyse ne constitue pas un audit RSE d'Isogeo. Elle vise plutôt à
mettre en évidence les liens qui peuvent être établis entre les activités de
l'entreprise, son organisation et certains principes associés à une démarche
de responsabilité sociétale.

\medskip

\noindent
\textbf{Enjeux économiques et gouvernance des données}

\medskip

La principale contribution de la solution Isogeo concerne la gouvernance et la
valorisation du patrimoine de données des organisations. La mise en place d'un
catalogue permet de mieux identifier les ressources existantes, leurs
producteurs, leurs conditions d'utilisation et leur niveau de documentation.

Cette connaissance favorise la réutilisation des données disponibles et limite
le risque de produire plusieurs fois une information déjà existante. Elle
contribue également à améliorer la traçabilité des données et à sécuriser les
processus métiers qui en dépendent.

Pour les organismes publics, une meilleure documentation facilite également la
diffusion et l'ouverture des données. Elle peut ainsi contribuer à renforcer la
transparence de l'action publique et l'accès des citoyens aux informations
produites par les administrations.
\footnote{%
    \href{https://www.isogeo.com/actualites/simplifiez-louverture-de-vos-donnees-geospatiales-avec-loffre-passerelle-donnees-quebec}
    {Isogeo, exemple de catalogage et de publication de données ouvertes,
    consulté en juillet 2026.}%
}

\medskip

\noindent
\textbf{Enjeux environnementaux}

\medskip

Les activités numériques ont un impact environnemental lié notamment à
l'utilisation des infrastructures informatiques, au stockage des données et à
la consommation de ressources de calcul. La gestion plus rigoureuse des données
peut contribuer indirectement à limiter certains traitements inutiles ou
redondants.

Le fonctionnement de Scan illustre cette logique. Une signature est calculée à
partir du contenu des données afin de détecter leur modification. Lorsqu'une
ressource n'a pas changé depuis le scan précédent, il n'est pas nécessaire de
réaliser à nouveau l'intégralité de sa documentation. Ce mécanisme permet
d'optimiser le temps de traitement et d'éviter certaines opérations
inutiles.
\footnote{%
    \href{https://www.isogeo.com/actualites/r-d-isogeo---ensg-detection-de-la-modification-des-tables-geographiques}
    {Isogeo, « Détection de la modification des tables géographiques »,
    septembre 2022.}%
}

La migration de traitements FME vers des composants Python autonomes peut
également participer à une rationalisation de l'architecture technique en
réduisant le nombre de dépendances nécessaires au fonctionnement du produit.
Il convient toutefois de rester prudent : aucun gain environnemental précis ne
peut être affirmé sans mesure comparative de la consommation de ressources
avant et après la migration.

\medskip

\noindent
\textbf{Enjeux sociaux et qualité du travail}

\medskip

La documentation et le partage des connaissances occupent une place importante
dans le fonctionnement de l'équipe. Les revues de code, les échanges techniques
et la rédaction de documentation contribuent à rendre les développements plus
compréhensibles et à limiter la dépendance à une seule personne.

L'accueil d'étudiants en stage ou en alternance participe également à la
transmission des compétences dans les domaines du développement logiciel, de la
géomatique et de la gouvernance des données.

Enfin, les produits développés par Isogeo cherchent à rendre les données
géographiques plus facilement accessibles à leurs utilisateurs. En facilitant
la recherche et la compréhension des données, la solution contribue à réduire
les obstacles techniques qui peuvent limiter leur exploitation.

% =============================================================================
% 1.2 PRÉSENTATION DU SUJET D'ALTERNANCE
% =============================================================================

\section{Présentation du sujet}
\label{sec:sujet_alternance}

% -----------------------------------------------------------------------------
% 1.2.1 Contexte technique du produit Scan
% -----------------------------------------------------------------------------

\subsection{Contexte technique du produit Scan}
\label{subsec:contexte_scan}

Scan est le composant de la solution Isogeo chargé d'explorer les sources de
données configurées par un utilisateur. Il identifie les ressources présentes,
en extrait différentes informations techniques, puis transmet les résultats à
la plateforme Isogeo afin de créer ou de mettre à jour les fiches de
métadonnées correspondantes.

Le service installé dans l'environnement du client est appelé
\textit{Isogeo Worker}. Il s'exécute sur une machine Windows disposant des
droits nécessaires pour consulter les sources ciblées. Les opérations réalisées
par Scan sont des opérations de lecture : elles n'ont pas pour objectif de
modifier les données des clients.
\footnote{%
    \href{https://help.isogeo.com/doc-scan-fme/support/exe/}
    {Documentation officielle de Scan, présentation du sous-programme
    d'extraction de métadonnées, consultée en juillet 2026.}%
}

Le traitement d'une source peut être résumé en trois opérations principales.

\medskip

\noindent
\textbf{Le listing}

\medskip

Le listing consiste à inventorier les ressources accessibles à partir d'un
point d'entrée. Selon le type de source, une ressource peut correspondre à une
table, une vue, une couche géographique, un fichier ou un jeu de données.

Cette étape permet à Scan de connaître la liste des éléments disponibles avant
de décider lesquels doivent être analysés plus en détail.

\medskip

\noindent
\textbf{La signature}

\medskip

La signature correspond à une empreinte calculée à partir du contenu d'une
ressource. Elle permet à Scan de reconnaître une donnée déjà rencontrée et de
détecter si son contenu a été modifié entre deux exécutions.

Cette information joue un rôle important dans la mise à jour du catalogue. Si
une donnée n'a pas changé, certains traitements peuvent être évités. Si sa
signature évolue, Scan peut au contraire déclencher une nouvelle extraction de
ses métadonnées.

\medskip

\noindent
\textbf{L'extraction des métadonnées}

\medskip

L'étape d'extraction, également appelée \textit{lookup} dans la documentation
technique de Scan, produit les informations décrivant la ressource analysée.

Selon la nature de la donnée, les informations extraites peuvent notamment
comprendre :

\begin{itemize}
    \item le nom et l'emplacement de la ressource ;
    \item son format ;
    \item le système de coordonnées ;
    \item le nombre d'entités ou d'enregistrements ;
    \item le type de géométrie ;
    \item les coordonnées de l'emprise spatiale ;
    \item une enveloppe convexe ;
    \item les noms, types, alias, tailles et précisions des champs ;
    \item certaines propriétés propres aux données raster ;
    \item différentes informations techniques utiles à la création de la fiche
          de métadonnées.
\end{itemize}

Les résultats sont structurés de manière à pouvoir être interprétés par les
autres composants de Scan. La nouvelle implémentation Python doit donc respecter
le contrat de données déjà utilisé par le produit, notamment la structure des
résultats, les paramètres d'exécution et les codes d'erreur.

\medskip

\noindent
\textbf{Diversité des sources prises en charge}

\medskip

Scan doit fonctionner dans des environnements techniques variés. Les données
peuvent être stockées dans :

\begin{itemize}
    \item une base PostgreSQL avec l'extension spatiale PostGIS ;
    \item une base Oracle ou Oracle Spatial ;
    \item une base Microsoft SQL Server ;
    \item une géodatabase d'entreprise Esri ;
    \item une géodatabase fichier Esri ;
    \item différents formats de fichiers vectoriels, raster ou tabulaires ;
    \item certains services géographiques.
\end{itemize}

Cette diversité constitue l'une des principales difficultés du produit. Chaque
source possède ses propres méthodes de connexion, ses types de données, ses
fonctions spatiales et ses comportements particuliers. Une opération simple,
comme déterminer le nombre d'entités ou calculer une enveloppe convexe, peut
ainsi nécessiter une implémentation différente selon la technologie utilisée.

% -----------------------------------------------------------------------------
% 1.2.2 Problématique de migration des traitements FME
% -----------------------------------------------------------------------------

\subsection{Problématique de migration des traitements FME}
\label{subsec:problematique_fme}

FME est une plateforme spécialisée dans l'intégration et la transformation de
données. Elle propose de nombreux lecteurs, écrivains et composants de
traitement adaptés aux données géographiques. Elle a donc historiquement occupé
une place importante dans le fonctionnement de Scan.

Plusieurs traitements étaient décrits au moyen de workbenches FME. Ces chaînes
graphiques permettaient notamment de se connecter aux sources, de parcourir les
ressources disponibles, de calculer certaines informations géométriques et de
produire les résultats attendus par Scan.

Le choix d'engager une migration ne remet pas en cause les capacités
fonctionnelles de FME. Il répond avant tout à de nouveaux besoins
d'industrialisation et d'autonomie du produit.

\medskip

\noindent
\textbf{Dépendance à un logiciel externe}

\medskip

L'exécution des workbenches nécessite la présence d'une installation FME
compatible ainsi que d'une licence disponible. Le fonctionnement de Scan dépend
donc d'un logiciel tiers qui doit être installé, configuré et maintenu dans
l'environnement du client.

Cette dépendance augmente le nombre de prérequis techniques nécessaires au
déploiement du produit. Elle peut également compliquer les opérations de support
lorsqu'un problème dépend de la version de FME, de l'activation de la licence ou
d'un composant particulier de l'installation.

La documentation de support de Scan indique notamment que la disponibilité de
la licence FME et la compatibilité de l'environnement font partie des éléments
à contrôler lors du diagnostic d'une exécution.
\footnote{%
    \href{https://help.isogeo.com/doc-scan-fme/support/support/}
    {Documentation officielle de support de Scan, consultée en juillet 2026.}%
}

\medskip

\noindent
\textbf{Maintenance des traitements}

\medskip

Les workbenches FME offrent une représentation graphique des traitements, mais
leur maintenance peut devenir complexe lorsque le nombre de scénarios, de
sources et de cas particuliers augmente.

Certaines logiques peuvent être réparties entre plusieurs composants ou
plusieurs fichiers. La comparaison des versions, la factorisation du code et la
mise en place de tests automatisés sont également différentes de celles
habituellement utilisées dans un projet logiciel développé en Python.

Une implémentation Python permet de regrouper les comportements communs dans des
classes et des fonctions réutilisables. Elle facilite également l'utilisation
d'outils de contrôle de qualité, de tests unitaires, de tests d'intégration et
de gestion de versions.

\medskip

\noindent
\textbf{Intégration dans l'architecture de Scan}

\medskip

Le remplacement des workbenches par des composants Python permet de rapprocher
les traitements de données du reste du code applicatif. Les paramètres, les
journaux d'exécution, les erreurs et les résultats peuvent être gérés de manière
plus homogène.

La solution peut également être distribuée sous la forme d'un exécutable
Windows autonome. L'objectif est que le client puisse exécuter les traitements
sans avoir à installer Python ni à gérer manuellement l'ensemble des
bibliothèques utilisées.

\medskip

\noindent
\textbf{Besoin de continuité fonctionnelle}

\medskip

La migration ne peut toutefois pas se limiter à produire des résultats
approximativement similaires. Les nouveaux traitements doivent conserver le
comportement attendu par les autres composants de Scan.

Ils doivent notamment respecter :

\begin{itemize}
    \item les paramètres de connexion existants ;
    \item la structure des données retournées ;
    \item les conventions de nommage ;
    \item les règles de calcul des métadonnées ;
    \item les codes d'erreur attendus ;
    \item les performances nécessaires au traitement de volumes importants ;
    \item les particularités des différentes sources.
\end{itemize}

La difficulté principale réside donc dans la reproduction fidèle d'un
comportement existant tout en modifiant profondément la technologie utilisée
pour l'obtenir.

Ce chantier s'inscrit dans une évolution plus générale du produit. Isogeo a
annoncé en 2026 l'avancement du Scan sans FME, avec pour objectif de supprimer
progressivement FME des prérequis nécessaires à la plateforme.
\footnote{%
    \href{https://www.isogeo.com/actualites/2-journee-communaute-isogeo-retours-terrain-nouveautes-produit-et-ia-geospatiale}
    {Isogeo, « 2\ieme{} Journée Communauté Isogeo », juin 2026.}%
}

% -----------------------------------------------------------------------------
% 1.2.3 Reformulation du sujet
% -----------------------------------------------------------------------------

\subsection{Reformulation du sujet}
\label{subsec:reformulation_sujet}

Le sujet de cette alternance peut être reformulé de la manière suivante :

\begin{quote}
    \itshape
    Concevoir, développer et valider une solution Python capable de remplacer
    progressivement les traitements FME utilisés par Scan pour inventorier et
    documenter différentes sources de données, puis préparer son exécution
    autonome dans un environnement Windows.
\end{quote}

La mission ne consistait donc pas à développer un nouveau produit indépendant,
mais à faire évoluer un composant existant en conservant sa compatibilité avec
le reste de l'écosystème Isogeo.

\medskip

\noindent
\textbf{Périmètre de la mission}

\medskip

Le périmètre des travaux comprenait principalement :

\begin{itemize}
    \item l'étude des workbenches FME existants ;
    \item l'identification des traitements à migrer ;
    \item la compréhension des formats d'entrée et de sortie ;
    \item le développement de composants Python dédiés aux différentes sources ;
    \item l'utilisation de GDAL/OGR pour les traitements de données
          géographiques ;
    \item l'utilisation de bibliothèques ou d'API spécifiques lorsque
          GDAL/OGR ne permettait pas de reproduire le comportement attendu ;
    \item la gestion des paramètres de connexion ;
    \item l'extraction des attributs et des propriétés géométriques ;
    \item la gestion des erreurs et des données invalides ;
    \item la comparaison des résultats Python avec ceux de FME ;
    \item la construction de jeux de tests de régression ;
    \item la préparation du packaging de l'application avec PyInstaller ;
    \item la production d'un exécutable Windows destiné à fonctionner sans
          installation locale de Python.
\end{itemize}

Les principales sources étudiées comprenaient notamment PostgreSQL/PostGIS,
Oracle, Microsoft SQL Server, les géodatabases Esri ainsi que plusieurs formats
de données vectorielles et raster.

\medskip

\noindent
\textbf{Limites du travail}

\medskip

Le projet s'inscrit dans une migration progressive. Il n'avait donc pas pour
objectif de réécrire simultanément l'ensemble du produit Scan ni de prendre en
charge immédiatement tous les formats supportés historiquement par FME.

Le travail portait principalement sur les traitements de lecture, d'inventaire,
de calcul et d'extraction des métadonnées. Les interfaces web de Scan, la
plateforme centrale Isogeo et les fonctionnalités de gestion des catalogues ne
faisaient pas partie du périmètre principal de la mission.

Certaines limitations dépendaient également des technologies utilisées. Les
géodatabases Esri nécessitent, par exemple, la présence de composants et de
licences propres à l'écosystème ArcGIS. De même, les comportements des systèmes
de gestion de bases de données peuvent varier en fonction de leur version, de
leur configuration et des droits accordés à l'utilisateur.

Enfin, les traitements FME existants ont été conservés comme référence pendant
la phase de migration. Leur maintien permettait de comparer les résultats,
d'identifier les écarts et de vérifier que la nouvelle solution respectait le
comportement historique attendu par Scan.

La réussite de la mission dépendait donc autant de la qualité du développement
Python que de la mise en place d'une méthode de validation suffisamment fiable
pour démontrer l'équivalence des traitements.